\documentclass[10pt,a4paper]{amsart}
\usepackage[margin=19mm]{geometry}
\usepackage{amsmath,amssymb,mathtools}
\usepackage[scr=rsfs,cal=euler]{mathalfa}
\usepackage{xfrac}
\usepackage[unicode,hidelinks,bookmarks=false]{hyperref}
\newcommand{\ie}{i.e.\ }
\newcommand{\cf}{cf.\ }
\newcommand{\resp}{respectively\ }
\newcommand{\Subsection}{\subsection}
\providecommand{\nobreakdash}{\nobreak\hbox{-}}
\setlength{\emergencystretch}{2em}
\multlinegap0pt
\usepackage{bm}
\usepackage{bbm}
\usepackage{dsfont}
\usepackage{xspace}
\usepackage{cancel}
\usepackage{mathtools}
\usepackage{amsthm}
\makeatletter
\@ifundefined{definition}{\newtheorem{definition}{Definition}}{}
\@ifundefined{lemma}{\newtheorem{lemma}{Lemma}}{}
\makeatother
\DeclareMathOperator{\even}{even}
\DeclareMathOperator{\odd}{odd}
\newcommand{\QQ}{\mathbb{Q}}
\newcommand{\ZZ}{\mathbb{Z}}
\newcommand{\xx}{\mathbf{x}}
\newcommand{\yy}{\mathbf{y}}
\renewcommand{\geq}{\geqslant}
\renewcommand{\leq}{\leqslant}
\title[Random Diophantine equations of large degree]{Random Diophantine equations of large degree}
\author{Tim Browning, Will Sawin}
\date{Source: arXiv:2510.26191v1 (CC BY 4.0)}
\begin{document}
\maketitle

\noindent\emph{Source and licence.} Random Diophantine equations of large degree. Version arXiv:2510.26191v1 (CC BY 4.0), distributed under the Creative Commons Attribution 4.0 International licence, as recorded in the arXiv licence metadata for this exact version and shown on the abstract page https://arxiv.org/abs/2510.26191v1. Full rights notice: licenses/arxiv-2510.26191-CC-BY-4.0.txt.\par\smallskip

\noindent\textit{Editorial scope.} Counted unit: the Lemma whose source label is \texttt{Lemma second moment} in the article (source0.tex lines 584--596, source label \texttt{Lemma second moment}), delivered together with the article's own complete original proof of it (source0.tex lines 598--1065, one \texttt{proof} environment of 468 lines). No mathematics is written, repaired or re-derived: statement and proof are the article's own text line for line. Required context retained uncounted: Cardinality Bdn (lines 99--102); Definition c (lines 221--224); Estimate Stirling binom (lines 317--320); Upper bound Stirling binom (lines 322--326); Identity VDM (lines 335--338); Identity VDM' (lines 341--344); Identity VDM2 (lines 347--350); Lemma coefficients (lines 390--396); lem:prob (lines 466--472); Cardinality Unl (lines 508--511); def-f (lines 519--521); Lemma first moment (lines 525--531); eq:3.2 (lines 554--563); def-f' (lines 571--580). Every definition, display and label of those lines is retained unchanged. Omitted from this package and not redistributed: 1--98, 103--220, 225--316, 321--321, 327--334, 339--340, 345--346, 351--389, 397--465, 473--507, 512--518, 522--524, 532--553, 564--570, 581--583, 597--597, 1066--1375. Nothing inside a retained excerpt is omitted or altered: every retained body is delivered exactly as the article's lines are written, so no omission receipt of any kind accompanies this package. Citations: the retained excerpts carry no citation command at all, so no bibliography entry of the article is transcribed and no reference is redistributed. Reference adaptation: none was needed and none was made; every reference inside the delivered unit resolves inside it, so no unresolved reference or citation is produced. Printed slips kept literal: no printed word, symbol, hypothesis or number of the counted statement or of its proof was altered. Layout and class: the article's own class is \texttt{amsart} with packages including amsmath, amsthm, amssymb, amsfonts, mathrsfs, geometry, color, hyperref; the wrapper is the lane's pinned \texttt{amsart} editorial wrapper at 10pt on A4, which restates the theorem-like environments the retained text needs and adds the article's own macro definitions that the retained text needs (\texttt{\textbackslash QQ}, \texttt{\textbackslash ZZ}, \texttt{\textbackslash even}, \texttt{\textbackslash geq}, \texttt{\textbackslash leq}, \texttt{\textbackslash odd}, \texttt{\textbackslash xx}, \texttt{\textbackslash yy}), with extra wrapper packages (bm, bbm, dsfont, xspace, cancel, mathtools). The delivered numbering is the unit's own. Assets: no retained span contains a figure environment, a graphics-inclusion command, a PDF-page inclusion, a table or a TikZ picture; no image file is copied into this package, the package declares no asset dependency and no image file is redistributed with it. Licence: version arXiv:2510.26191v1 of this article is distributed under the Creative Commons Attribution 4.0 International licence, as recorded in the arXiv licence metadata for this exact version and shown on the abstract page https://arxiv.org/abs/2510.26191v1. A full rights notice is delivered at licenses/arxiv-2510.26191-CC-BY-4.0.txt. Characters: every retained excerpt is pure ASCII, so the delivered engine, XeLaTeX with its default Latin Modern text font, reports no missing character.
\begin{equation}
\label{Cardinality Bdn}
\# \mathbb{B}_{d,n} = 2^{N_{d,n}-1}.
\end{equation}

\begin{equation}
\label{Definition c}
c_{d,\ell}(j) = \sum_{\substack{k=0 \\ 2 \nmid k}}^d \binom{k+j-1}{k} \binom{d-k+\ell-j}{d-k}.
\end{equation}

\begin{equation}
\label{Estimate Stirling binom}
\binom{2m}{m} = \frac{2^{2m}}{\pi^{1/2} m^{1/2}} \left(1 + O \left( \frac1{m} \right) \right).
\end{equation}

\begin{equation}
\label{Upper bound Stirling binom}
\binom{2m}{m} \leq \frac{2^{2m}}{\pi^{1/2} m^{1/2}} \quad \text{ and } \quad 
\binom{m}{\lfloor m/2 \rfloor} \leq \frac{2^{m+1/2}}{\pi^{1/2} m^{1/2}},
\end{equation}

\begin{equation}
\label{Identity VDM}
\binom{r_1+r_2}{q} = \sum_{\alpha_1=0}^{r_1} \binom{r_1}{\alpha_1} \binom{r_2}{q-\alpha_1},
\end{equation}

\begin{equation}
\label{Identity VDM'}
\binom{r_1+r_2}{r_2-q} = \sum_{\alpha_1=0}^{r_1} \binom{r_1}{\alpha_1} \binom{r_2}{q+\alpha_1},
\end{equation}

\begin{equation}
\label{Identity VDM2}
\binom{r_1+r_2+r_3}{q} = \sum_{\alpha_1=0}^{r_1} \ \sum_{\alpha_2=0}^{r_2} \binom{r_1}{\alpha_1} \binom{r_2}{\alpha_2} \binom{r_3}{q-\alpha_1-\alpha_2},
\end{equation}

\begin{lemma}
\label{Lemma coefficients}
Let $1\leq j\leq \ell \leq d$. Then we have 
\begin{equation*}
\left|1- \frac{2c_{d,\ell}(j)}{N_{d,\ell}}\right|  \leq \left( \frac{1}{2} \right)^{\min \{ j, \ell+1-j \}}.
\end{equation*}
\end{lemma}

\begin{lemma}\label{lem:prob}
Let $r,s\geq 0$ be integers. Then 
$$
\sum_{j=0}^r\binom{r}{j}\binom{s}{j} 2^{-j}\leq \left(\frac{3}{4}\right)^\mu \binom{r+s}{r},
$$
where $\mu=rs/(r+s).$
\end{lemma}

\begin{equation}
\label{Cardinality Unl}
\# U_{n,\ell} = 2^{\ell} \binom{n+1}{\ell+1}.
\end{equation}

\begin{definition}\label{def-f}
Let $n\geq 1$. Given $f\in \ZZ[x_0,\dots,x_n]$, we let $\nu(f)$ denote the number of coefficients of $f$ equal to $-1$.
\end{definition}

\begin{lemma}
\label{Lemma first moment}
Let $d, n \geq 2$ and let $\ell \in \{1, \dots, n\}$ be such that $N_{d,\ell}$ is even. We have
\begin{equation*}
M_1(d,n;\ell) = \frac1{2^{N_{d,\ell}-\ell}}  \binom{n+1}{\ell+1}\binom{N_{d,\ell}}{N_{d,\ell}/2}.
\end{equation*}
\end{lemma}

\begin{equation}\label{eq:3.2}
\begin{split}
\sum_{\substack{V \in \mathbb{B}_{d,n} \\ x \in V(\mathbb{Q})}} 1
&=
2^{N_{d,n}-N_{d,\ell}}
\#\left\{
V \in \mathbb{B}_{d,\ell}: \nu(f_V)=N_{d,\ell}/2\right\}\\
&=2^{N_{d,n}-N_{d,\ell}-1} \binom{N_{d,\ell}}{N_{d,\ell}/2},
\end{split}
\end{equation}

\begin{definition}\label{def-f'}
Let $n\geq 1$
and let $j\in \{1,\dots,n\}.$
Given $f\in \ZZ[x_0,\dots,x_n]$, we define  $\even_j(f)$ to be the part of $f$ that is 
even with  respect to the first $j$ variables. Similarly, we define 
$\odd_j(f)$ to be the part of $f$ that is  odd  with respect to the first $j$ variables. 
Next, let 
$a\in \{1,\dots,n\}$ and  $b\in \{0,\dots,n+1-a\}$.
We also set $R_a^b(f)$ to be the sum of the monomials of $f$ involving only the first $a$ and the final $b$ variables.
\end{definition}

\begin{lemma}
\label{Lemma second moment}
Let $d, n \geq 2$ and let $\ell \in \{1, \dots, \min\{d,n\}\}$ be such that $N_{d,\ell}$ is even. Then 
\begin{equation*}
\frac{M_2(d,n;\ell)}{M_1(d,n;\ell)^2} \leq 
1 + \boldsymbol{1}_{\ell=n} + O \left( \left(\frac{3}{4}\right)^{\min\{\ell,\mu_\ell\}}+\frac1{\min\{d,n\}}+
 \frac1{M_1(d,n;\ell)}\right) ,
\end{equation*}
where
\begin{equation}\label{eq:soil}
\mu_\ell=\frac{(n-\ell)(\ell+1)}{n+1}.
\end{equation}
\end{lemma}

\begin{proof}
We start by noting that
\begin{equation*}
M_2(d,n;\ell) = \frac1{\# \mathbb{B}_{d,n}} \sum_{x, y \in U_{n,\ell}} \sum_{\substack{V \in \mathbb{B}_{d,n} \\ x, y \in V(\mathbb{Q})}} 1.
\end{equation*}
We set $m_{n,\ell} = \max\{2\ell-n,0\}$ and for $i \in \{m_{n,\ell}, \dots, \ell\}$ and $j \in \{0, \dots, i\}$, we let $T_{n,\ell}^{(i,j)}$ denote the set of pairs of elements $x, y \in U_{n,\ell}$ which have exactly $i+1$ nonzero coordinates in common, and are such that among these $i+1$ coordinates exactly $j$ are different when the final nonzero coordinate shared by $x$ and $y$ is viewed as being $1$. Moreover, we let $T_{n,\ell}$ denote the set of pairs of elements of $ U_{n,\ell}$ which do not have any nonzero coordinates in common. We notice that
\begin{equation*}
U_{n,\ell} \times U_{n,\ell} = \left( \bigsqcup_{i=m_{n,\ell}}^{\ell} \ \bigsqcup_{j=0}^i T_{n,\ell}^{(i,j)} \right) \sqcup T_{n,\ell},
\end{equation*}
which allows us to write
\begin{equation}
\label{Equality M2}
M_2(d,n;\ell) = M_1(d,n;\ell) + \Sigma^{(1)}(d,n;\ell)+\Sigma^{(2)}(d,n;\ell)+\Sigma^{(3)}(d,n;\ell),
\end{equation}
where
\begin{equation*}
\Sigma^{(1)}(d,n;\ell) = \frac1{\# \mathbb{B}_{d,n}} \sum_{j=1}^{\ell} \ \sum_{(x, y) \in T_{n,\ell}^{(\ell,j)}} \sum_{\substack{V \in \mathbb{B}_{d,n} \\ x, y \in V(\mathbb{Q})}} 1,
\end{equation*}
and
\begin{equation*}
\Sigma^{(2)}(d,n;\ell) = \frac1{\# \mathbb{B}_{d,n}} \sum_{i=m_{n,\ell}}^{\ell-1} \ \sum_{j=0}^i \ \sum_{(x, y) \in T_{n,\ell}^{(i,j)}} \sum_{\substack{V \in \mathbb{B}_{d,n} \\ x, y \in V(\mathbb{Q})}} 1,
\end{equation*}
and
\begin{equation*}
\Sigma^{(3)}(d,n;\ell) = \frac1{\# \mathbb{B}_{d,n}} \sum_{(x, y) \in T_{n,\ell}} \sum_{\substack{V \in \mathbb{B}_{d,n} \\ x, y \in V(\mathbb{Q})}} 1.
\end{equation*}
We note  that $\Sigma^{(2)}(d,n;\ell)=0$ if $\ell=n$ and $\Sigma^{(3)}(d,n;\ell)=0$ if $\ell \geq n/2$.
The most difficult sum to estimate is  
$\Sigma^{(2)}(d,n;\ell)$. While it ought to be possible to handle the three sums simultaneously, 
by allowing $i$ to also run over $-1$ and $\ell$ in the definition of $\Sigma^{(2)}(d,n;\ell)$, the treatment is  made easier through the restriction $i\leq \ell-1$, since we can then exploit a non-trivial upper bound for $N_{d,i}/N_{d,\ell}.$ 



Let $i\in \{m_{n,\ell},\dots,\ell\}$ and let $j\in \{0,\dots,i\}$. If $(x,y)\in T_{n,\ell}^{(i,j)}$ then 
$$
\sum_{\substack{V \in \mathbb{B}_{d,n} \\ x, y \in V(\mathbb{Q})}} 1=
2^{N_{d,n}-N_{d,2\ell-i}} 
\sum_{\substack{V \in \mathbb{B}_{d,2\ell-i} \\ \tilde x, \tilde y \in V(\mathbb{Q})}} 1,
$$
where $\tilde x, \tilde y\in \mathbb{P}^{2\ell-i}(\QQ)$ are obtained from $x$ and $y$ by removing their $n-2\ell+i$ common zeros. There exists an obvious bijection 
$\iota:  \mathbb{B}_{d,2\ell-i}\to \mathbb{B}_{d,2\ell-i}$ such that 
$$
\iota\left(
\left\{ 
V \in \mathbb{B}_{d,2\ell-i}: \tilde x, \tilde y \in V(\mathbb{Q})\right\}\right)=
\left\{ 
V \in \mathbb{B}_{d,2\ell-i} : f_V(\xx_{i,j})=f_V(\yy_{i,j})=0\right\},
$$
where
$$
\xx_{(i,j)} = (\underbrace{\strut 1 , \ldots , 1}_{j} , \underbrace{\strut 1 , \dots , 1}_{i+1-j} , 
\underbrace{\strut 1 , \dots , 1}_{\ell-i} , \underbrace{\strut 0 , \dots , 0}_{\ell-i})
$$
and 
$$
\yy_{(i,j)} = (\underbrace{\strut -1 , \dots , -1}_{j} , \underbrace{\strut 1 , \dots , 1}_{i+1-j} , 
\underbrace{\strut 0 , \dots , 0}_{\ell-i} , \underbrace{\strut 1 , \dots , 1}_{\ell-i}).
$$
Hence for 
 $i\in \{m_{n,\ell},\dots,\ell\}$ and $j\in \{0,\dots,i\}$, we have 
\begin{equation}\label{eq:**}
\sum_{\substack{V \in \mathbb{B}_{d,n} \\ x, y \in V(\mathbb{Q})}} 1=
2^{N_{d,n}-N_{d,2\ell-i}} 
\#\left\{V \in \mathbb{B}_{d,2\ell-i} :
f_V(\xx_{i,j})=f_V(\yy_{i,j})=0\right\},
\end{equation}
for  any $(x,y)\in T_{n,\ell}^{(i,j)}$.

Next, we note that for any $x \in U_{n,\ell}$ and $i \in \{m_{n,\ell}, \dots, \ell\}$, $j \in \{0, \dots, i\}$, we have 
\begin{equation*}
\# \left \{ y \in U_{n,\ell} : (x, y) \in T_{n,\ell}^{(i,j)} \right\} = 2^{\ell-i} \binom{n-\ell}{\ell-i} \binom{\ell+1}{i+1} \binom{i}{j},
\end{equation*}
and
\begin{equation*}
\# \left \{ y \in U_{n,\ell} : (x, y) \in T_{n,\ell} \right\} = 2^{\ell} \binom{n-\ell}{\ell+1}.
\end{equation*}
Recalling the equality \eqref{Cardinality Unl}, we thus deduce that
\begin{equation}
\label{Cardinality Tnlij}
\# T_{n,\ell}^{(i,j)} = 2^{2\ell-i} \binom{n+1}{\ell+1} \binom{n-\ell}{\ell-i} \binom{\ell+1}{i+1} \binom{i}{j},
\end{equation}
and
\begin{equation}
\label{Cardinality Tnl}
\# T_{n,\ell} = 2^{2 \ell} \binom{n+1}{\ell+1} \binom{n-\ell}{\ell+1}.
\end{equation}

\subsubsection*{Analysis of $\Sigma^{(1)}(d,n;\ell)$}


This concerns the case $i=\ell$, so that $2\ell-i=\ell$. 
Given $V\in 
 \mathbb{B}_{d,\ell}$ and $j\in \{1,\dots,\ell\}$, we note that
\begin{align*}
f_V(\xx_{i,j})
&=N_{d,\ell}-2\nu(f_V)\\
&=N_{d,\ell}-2\nu(\even_j(f_V))-2\nu(\odd_j(f_V)),
\end{align*}
in the notation of Definition \ref{def-f'}.
Recall that 
 $c_{d,\ell}(j)$ is defined to be the number of monomials of degree $d$ which are odd with respect to $j$ given variables among $\ell+1$, and that it satisfies  \eqref{Definition c}.
Then we find similarly that
 \begin{align*}
f_V(\yy_{i,j})
&=N_{d,\ell}-
c_{d,\ell}(j)-
2\nu(\even_j(f_V))-\left(c_{d,\ell}(j)-2\nu(\odd_j(f_V))\right)\\
&=N_{d,\ell}-
2c_{d,\ell}(j)-
2\nu(\even_j(f_V))+2\nu(\odd_j(f_V)).
\end{align*}
Hence it follows from \eqref{eq:**} that 
$$
\sum_{\substack{V \in \mathbb{B}_{d,n} \\ x, y \in V(\mathbb{Q})}} 1=
2^{N_{d,n}-N_{d,\ell}} 
\#\left\{V \in \mathbb{B}_{d,\ell} : 
\begin{array}{l}
\nu(\even_j(f_V))=(N_{d,\ell}-c_{d,\ell}(j))/2,\\
\nu(\odd_j(f_V))=c_{d,\ell}(j)/2
\end{array}\right\}.
$$
Since the integer $N_{d,\ell}$ is assumed to be even, we deduce that 
\begin{equation*}
\sum_{\substack{V \in \mathbb{B}_{d,n} \\ x, y \in V(\mathbb{Q})}} 1 =
\begin{cases}
0 & \textrm{if } 2 \nmid c_{d,\ell}(j), \\
\displaystyle{2^{N_{d,n}-N_{d,\ell}-1} \binom{N_{d,\ell}-c_{d,\ell}(j)}{(N_{d,\ell}-c_{d,\ell}(j))/2} \binom{c_{d,\ell}(j)}{c_{d,\ell}(j)/2}} & \textrm{if } 2 \mid c_{d,\ell}(j),
\end{cases}
\end{equation*}
for any  $j \in \{1, \dots, \ell\}$ and any $(x, y) \in T_{n,\ell}^{(\ell,j)}$.
Noting that $0<c_{d,j}(\ell)<N_{d,\ell}$, we may
apply the upper bound \eqref{Upper bound Stirling binom}
 twice to  derive the inequality
\begin{equation*}
\sum_{\substack{V \in \mathbb{B}_{d,n} \\ x, y \in V(\mathbb{Q})}} 1 \leq \frac{2^{N_{d,n}+1}}{\pi N_{d,\ell} } \left( 1-\left|
 1-\frac{2c_{d,\ell}(j)}{N_{d,\ell}}\right|^2\right)^{-1/2}.
\end{equation*}
Since  $\ell \leq d$ by assumption, we 
 see that we are in position to apply Lemma~\ref{Lemma coefficients}. 
Thus 
 $
 \left|
 1-\frac{2c_{d,\ell}(j)}{N_{d,\ell}}\right|\leq 0.5,
 $
and we may deduce that
 \begin{equation*}
\sum_{\substack{V \in \mathbb{B}_{d,n} \\ x, y \in V(\mathbb{Q})}} 1 \leq \frac{2^{N_{d,n}+1}}{\pi N_{d,\ell} } \left( 1+O\left(\left|
 1-\frac{2c_{d,\ell}(j)}{N_{d,\ell}}\right|^2\right)\right).
\end{equation*}
 But then a further  
application of   Lemma~\ref{Lemma coefficients} gives
\begin{equation*}
\sum_{\substack{V \in \mathbb{B}_{d,n} \\ x, y \in V(\mathbb{Q})}} 1 \leq \frac{2^{N_{d,n}+1}}{\pi N_{d,\ell}} \left( 1 + O \left( \left(\frac{1}{4} \right)^{\min \{ j, \ell-j \}} \right) \right).
\end{equation*}
Appealing to the equalities \eqref{Cardinality Bdn} and \eqref{Cardinality Tnlij}, we deduce that
\begin{equation*}
\Sigma^{(1)}(d,n;\ell) \leq \frac{2^{\ell+2}}{\pi N_{d,\ell}} \binom{n+1}{\ell+1} \sum_{j=1}^{\ell} \binom{\ell}{j} \left( 1 + O \left( \left(\frac{1}{4} \right)^{\min \{ j, \ell-j \}} \right) \right).
\end{equation*}
But for any $A\geq 1$ we have 
$$
\sum_{a=0}^A \binom{A}{a}=2^A \quad \text{ and } \quad 
\sum_{a=0}^A \binom{A}{a}\left(\frac{1}{4}\right)^a=\left(\frac{5}{4}\right)^A.
$$
Thus 
\begin{equation*}
\sum_{j=1}^{\ell} \binom{\ell}{j} \left( 1 + O \left( \left(\frac{1}{4} \right)^{\min \{ j, \ell-j \}} \right) \right)=
2^\ell \left(1+ O\left(\left(\frac{5}{8}\right)^\ell\right)\right),
\end{equation*}
which yields  in particular that
\begin{align*}
\Sigma^{(1)}(d,n;\ell) & \leq  \frac{2^{2\ell+2}}{\pi N_{d,\ell}} \binom{n+1}{\ell+1} \left( 1 + O\left(
\left(\frac{3}{4}\right)^\ell\right) \right).
\end{align*}
As a result, on appealing to the estimate \eqref{Estimate Stirling binom} and noting that $N_{d,\ell} \geq  d$, we eventually get
\begin{equation}
\label{Upper bound Sigma1}
\begin{split}
\Sigma^{(1)}(d,n;\ell) 
&\leq \frac1{2^{2N_{d,\ell}-2\ell-1}} \binom{N_{d,\ell}}{N_{d,\ell}/2}^2 \binom{n+1}{\ell+1} \left( 1 + O \left( \left(\frac{3}{4}\right)^\ell+\frac{1}{d}\right) \right)\\
&= M_1(d,n;\ell)^2 
\left( \frac{2}{\binom{n+1}{\ell+1}} + O \left( \left(\frac{3}{4}\right)^\ell+\frac{1}{d}\right) \right),
\end{split}
\end{equation}
where the second line follows from Lemma \ref{Lemma first moment}.

\subsubsection*{Analysis of $\Sigma^{(2)}(d,n;\ell)$}

This concerns the case $i\in \{m_{n,\ell},\dots,\ell-1\}$ and $j\in \{0,\dots,i\}$.
Our starting point is the observation that \eqref{eq:**} holds, 
for  any $(x,y)\in T_{n,\ell}^{(i,j)}$.
Then for any 
$V\in 
 \mathbb{B}_{d,2\ell-i}$ we have 
 \begin{align*}
f_V(\xx_{i,j})
&=N_{d,\ell}-2\nu(R_{\ell+1}^0(f_V))
=
N_{d,\ell}-2\nu(R_{i+1}^0(f_V))-2\nu(X_{\ell,i}(f_V)),
\end{align*}
in the notation of Definition \ref{def-f}, where
$X_{\ell,i}(f_V)=R_{\ell+1}^0(f_V)-R_{i+1}^0(f_V)$.
Note that the polynomial $X_{\ell,i}(f_V)$ only has $N_{d,\ell}-N_{d,i}$ coefficients. 
We conclude that $f_V(\xx_{i,j})=0$ if and only if 
\begin{equation}\label{eq:sun}
\nu(X_{\ell,i}(f_V))=\frac{N_{d,\ell}}{2}-
\nu\left( 
\even_j\left(R_{i+1}^0(f_V)\right)\right)
-\nu\left( 
\odd_j\left(R_{i+1}^0(f_V)\right)\right).
\end{equation}
Similarly, we find that 
 $$
f_V(\yy_{i,j})=N_{d,\ell}-c_{d,\ell}(j)-
2\nu\left( 
\even_j\left(R_{i+1}^{\ell-i}(f_V)\right)\right)
-\left(c_{d,\ell}(j)-
2\nu\left( 
\odd_j\left(R_{i+1}^{\ell-i}(f_V)\right)\right)\right),
$$
in the notation of \eqref{Definition c}.
Thus   $f_V(\yy_{i,j})=0$ if and only if 
\begin{equation}\label{eq:moon}
\begin{split}
\nu\left(\even_j\left( Y_{\ell,i}(f_V)\right)\right)
=~&\frac{N_{d,\ell}}{2}-
c_{d,\ell}(j)-
\nu\left( 
\even_j\left(R_{i+1}^0(f_V)\right)\right)\\
&+\nu\left( 
\odd_j\left(R_{i+1}^0(f_V)\right)\right)
+\nu\left( 
\odd_j\left(Y_{\ell,i}(f_V)\right)\right),
\end{split}
\end{equation}
where 
$Y_{\ell,i}(f_V)=R_{i+1}^{\ell-i}(f_V)-R_{i+1}^0(f_V)$.
Note that the polynomial $Y_{\ell,i}(f_V)$ only involves the 
$\ell+i$
variables
$x_{0},\dots,x_{i}, x_{\ell+1},\dots,x_{2\ell-i}$,
with at least one of the variables $x_{\ell+1},\dots,x_{2\ell-i}$ appearing.

It now follows from \eqref{eq:**} that 
$$
\sum_{\substack{V \in \mathbb{B}_{d,n} \\ x, y \in V(\mathbb{Q})}} 1=
2^{N_{d,n}-N_{d,2\ell-i}} 
\#\left\{V \in \mathbb{B}_{d,2\ell-i} : 
\text{ \eqref{eq:sun} and \eqref{eq:moon} hold}
\right\}.
$$
We put 
\begin{align*}
a&=\nu\left( 
\even_j\left(R_{i+1}^0(f_V)\right)\right)\in \left\{0,\dots,N_{d,i}-c_{d,i}(j)\right\},\\
b&=\nu\left( 
\odd_j\left(R_{i+1}^0(f_V)\right)\right)\in \left\{0,\dots,c_{d,i}(j)\right\},\\
c&=\nu\left( 
\odd_j\left(Y_{\ell,i}(f_V)\right)\right)\in \left\{0,\dots,c_{d,\ell}(j)-c_{d,i}(j)\right\}.
\end{align*}
The total number of positions for these $-1$ coefficients is 
$$
\binom{N_{d,i}-c_{d,i}(j)}{a}
\binom{c_{d,i}(j)}{b}
\binom{c_{d,\ell}(j)-c_{d,i}(j)}{c}.
$$
Moreover, it follows from \eqref{eq:sun} that the number of positions for $-1$ in $X_{\ell,i}(f_V)$
$$
\binom{N_{d,\ell}-N_{d,i}}{N_{d,\ell}/2-(a+b)},
$$
on recalling that  the integer $N_{d,\ell}$ is assumed to be even. Similarly, 
it follows from \eqref{eq:moon} that the number of positions for $-1$ in $Y_{\ell,i}(f_V)$ 
is 
$$
\binom{N_{d,\ell}-N_{d,i} -(c_{d,\ell}(j)-c_{d,i}(j))}{N_{d,\ell}/2-c_{d,\ell}(j)-a+b+c}.
$$
For fixed $a,b,c$, among the $N_{d,i}$ monomials in $x_0,\dots,x_i$ only, $a+b$ coefficients are $-1$, and the others are $+1$. Similarly, the signs are already accounted for among the 
 $N_{d,\ell}-N_{d,i}$
monomials in 
$R_{\ell+1}^0(f_V)-R_{i+1}^0(f_V)$ , and in the $N_{d,\ell}-N_{d,i}$ monomials in 
$R_{i+1}^{\ell-i}(f_V)-R_{i+1}^0(f_V)$. Hence the number of free coefficients is 
$$
N_{d,2\ell-i}-N_{d,i}-2(N_{d,\ell}-N_{d,i})=
N_{d,2\ell-i}-2N_{d,\ell}+N_{d,i}.
$$
In conclusion, for any $j \in \{0, \dots, i\}$ and $(x, y) \in T_{n,\ell}^{(i,j)}$, we deduce that 
\begin{align*}
\sum_{\substack{V \in \mathbb{B}_{d,n} \\ x, y \in V(\mathbb{Q})}} 1 & =  2^{N_{d,n}-2N_{d,\ell}+N_{d,i}-1} 
\hspace{-0.2cm}
\sum_{a=0}^{N_{d,i}-c_{d,i}(j)} \ \sum_{b=0}^{c_{d,i}(j)}\
\sum_{c=0}^{c_{d,\ell}(j)-c_{d,i}(j)}
\hspace{-0.2cm}
 \binom{N_{d,i}-c_{d,i}(j)}{a} \binom{c_{d,i}(j)}{b} \\
&\quad\times \binom{c_{d,\ell}(j)-c_{d,i}(j)}{c}
 \binom{N_{d,\ell}-N_{d,i}}{N_{d,\ell}/2-(a+b)} 
\binom{N_{d,\ell}-N_{d,i} -(c_{d,\ell}(j)-c_{d,i}(j))}{N_{d,\ell}/2-c_{d,\ell}(j)-a+b+c} .
\end{align*}
We apply  \eqref{Identity VDM'} to execute the sum over $c$, finding that 
\begin{align*}
\sum_{\substack{V \in \mathbb{B}_{d,n} \\ x, y \in V(\mathbb{Q})}} 1 =~ & \ 2^{N_{d,n}-2N_{d,\ell}+N_{d,i}-1} \sum_{a=0}^{N_{d,i}-c_{d,i}(j)} \ \sum_{b=0}^{c_{d,i}(j)}\
 \binom{N_{d,i}-c_{d,i}(j)}{a} \binom{c_{d,i}(j)}{b} \\
&\times
 \binom{N_{d,\ell}-N_{d,i}}{N_{d,\ell}/2-(a+b)} 
\binom{N_{d,\ell}-N_{d,i}}{N_{d,\ell}/2-\left(N_{d,i}-c_{d,i}(j)-a+b\right)} .
\end{align*}
Note that 
$$
\binom{N_{d,\ell}-N_{d,i}}{N_{d,\ell}/2-\alpha} 
\leq \binom{N_{d,\ell}-N_{d,i}}{
\lfloor (N_{d,\ell}-N_{d,i})/2\rfloor
},
$$
for any integer $\alpha\geq 0$.
On appealing to  the higher Vandermonde identity
\eqref{Identity VDM2}, we may conclude that 
\begin{align*}
\sum_{\substack{V \in \mathbb{B}_{d,n} \\ x, y \in V(\mathbb{Q})}} 1 \leq ~ & \ 2^{N_{d,n}-2N_{d,\ell}+N_{d,i}-1} 
\binom{N_{d,\ell}-N_{d,i}}{
\lfloor (N_{d,\ell}-N_{d,i})/2\rfloor
}
\binom{N_{d,\ell}}{
N_{d,\ell}/2},
\end{align*}
for each $i\in \{m_{n,\ell},\dots, \ell-1\}$ and  $j \in \{0, \dots, i\}$ and $(x, y) \in T_{n,\ell}^{(i,j)}$.

Observe that
$$
 \sum_{j=0}^i \binom{i}{j}=2^i.
$$
We now bring in the expression \eqref{Cardinality Tnlij} for the 
cardinality of $T_{n,\ell}^{(i,j)}$, together with \eqref{Cardinality Bdn}, in order to deduce that 
$$
\Sigma^{(2)}(d,n;\ell) \leq 
\frac{1}{2^{2N_{d,\ell}-2\ell}}\binom{N_{d,\ell}}{
N_{d,\ell}/2}^2\binom{n+1}{\ell+1}^2\Delta,
$$
where
$$
\Delta
=
\binom{N_{d,\ell}}{
N_{d,\ell}/2}^{-1}\binom{n+1}{\ell+1}^{-1}
 \sum_{i=m_{n,\ell}}^{\ell-1} 
 2^{N_{d,i}} 
 \binom{n-\ell}{\ell-i} \binom{\ell+1}{i+1} 
 \binom{N_{d,\ell}-N_{d,i}}{
\lfloor (N_{d,\ell}-N_{d,i})/2\rfloor
}.
$$
It clearly follows from Lemma \ref{Lemma first moment} that 
\begin{equation}\label{eq:leaf}
\Sigma^{(2)}(d,n;\ell) \leq 
M_1(d,n;\ell)^2
\Delta.
\end{equation}

We now turn to an analysis of $\Delta$, beginning with the upper bound
\begin{equation}\label{eq:root}
\Delta\leq 
\binom{n+1}{\ell+1}^{-1}
 \sum_{i=0}^{\ell-1} 
 \binom{n-\ell}{\ell-i} \binom{\ell+1}{i+1} \Delta_i,
\end{equation}
where
$$
\Delta_i= 2^{N_{d,i}} 
\binom{N_{d,\ell}}{
N_{d,\ell}/2}^{-1}
 \binom{N_{d,\ell}-N_{d,i}}{
\lfloor (N_{d,\ell}-N_{d,i})/2\rfloor
}.
$$
Combining the second bound in \eqref{Upper bound Stirling binom} with 
\eqref{Estimate Stirling binom}, we obtain  
$$
\Delta_i\leq \left(1-\frac{N_{d,i}}{N_{d,\ell}}\right)^{-1/2}\left(1+O\left(\frac{1}{N_{d,i}}\right)\right).
$$
But, on recalling that $i\leq \ell-1$,  we observe that 
$$
\frac{N_{d,i}}{N_{d,\ell}}=\prod_{j=i+1}^\ell \frac{j}{d+j}\leq \frac{1}{2^{\ell-i}},
$$
since $\frac{j}{d+j}\leq 1/2$ for $j\leq \ell\leq d$. Hence
$
\Delta_i\leq 1+O( 2^{-(\ell-i)})$, 
from which it follows that 
$$
\Delta\leq 
\binom{n+1}{\ell+1}^{-1}
 \sum_{i=0}^{\ell-1} 
 \binom{n-\ell}{\ell-i} \binom{\ell+1}{i+1} \left(1+O\left( \frac{1}{2^{\ell-i}}\right)\right)
$$
in \eqref{eq:root}.
On the one hand, we have 
\begin{align*}
 \sum_{i=0}^{\ell-1} \binom{n-\ell}{\ell-i} \binom{\ell+1}{i+1}
 &= 
 \sum_{i=1}^{\ell} \binom{n-\ell}{1+\ell-i} \binom{\ell+1}{i}\\
 &=
 \sum_{i=0}^{\ell+1} \binom{n-\ell}{1+\ell-i} \binom{\ell+1}{i}-1-\binom{n-\ell}{\ell+1}\\
&= \binom{n+1}{\ell+1}-1-\binom{n-\ell}{\ell+1},
\end{align*}
by the  Vandermonde identity \eqref{Identity VDM}.
On the other hand, we see that 
\begin{align*}
\sum_{i=0}^{\ell-1} 
 \binom{n-\ell}{\ell-i} \binom{\ell+1}{i+1} \frac{1}{2^{\ell-i}}
&=
\sum_{i=1}^{\ell} 
 \binom{n-\ell}{\ell+1-i} \binom{\ell+1}{i} \frac{1}{2^{\ell+1-i}}\\
 &\leq 
\sum_{i=0}^{\ell+1} 
 \binom{n-\ell}{\ell+1-i} \binom{\ell+1}{i} \frac{1}{2^{\ell+1-i}}\\
 &=
\sum_{j=0}^{\ell+1} 
 \binom{n-\ell}{j} \binom{\ell+1}{j} \frac{1}{2^{j}}.
\end{align*}
Thus, in the light of Lemma \ref{lem:prob}, we deduce that 
$$
\binom{n+1}{\ell+1}^{-1}
\sum_{i=0}^{\ell-1} 
 \binom{n-\ell}{\ell-i} \binom{\ell+1}{i+1} \frac{1}{2^{\ell-i}}
\leq 
\left(\frac{3}{4}\right)^{\mu_\ell},
$$
where $\mu_\ell$ is given by \eqref{eq:soil}.
Putting this together in \eqref{eq:leaf}, we deduce that 
\begin{equation}\label{eq:stalk}
\begin{split}
\Sigma^{(2)}(d,n;\ell)
\leq~& M_1(d,n;\ell)^2\left(1
-\frac{1}{\binom{n+1}{\ell+1}}
-\frac{\binom{n-\ell}{\ell+1}}{\binom{n+1}{\ell+1}}\right) +O\left( 
\left(\frac{3}{4}\right)^{\mu_\ell}
\right).
\end{split}
\end{equation}


\subsubsection*{Analysis of $\Sigma^{(3)}(d,n;\ell)$}

We remark that reasoning exactly as in the proof of  \eqref{eq:3.2}, we deduce that 
\begin{equation*}
\sum_{\substack{V \in \mathbb{B}_{d,n} \\ x, y \in V(\mathbb{Q})}} 1 = 2^{N_{d,n}-2N_{d,\ell}-1} \binom{N_{d,\ell}}{N_{d,\ell}/2}^2,
\end{equation*}
for any  $(x, y) \in T_{n,\ell}$.
Recalling  \eqref{Cardinality Bdn} and \eqref{Cardinality Tnl}, we thus see that
\begin{equation}
\label{Equality Sigma3}
\begin{split}
\Sigma^{(3)}(d,n;\ell) 
&= \frac1{2^{2N_{d,\ell}-2\ell}} \binom{N_{d,\ell}}{N_{d,\ell}/2}^2 \binom{n+1}{\ell+1} \binom{n-\ell}{\ell+1}\\
&= M_1(d,n;\ell)^2  \frac{\binom{n-\ell}{\ell+1}}{\binom{n+1}{\ell+1}},
\end{split}
\end{equation}
by Lemma~\ref{Lemma first moment}.

\subsubsection*{Conclusion of the argument}

Noting that 
$\binom{n+1}{\ell+1} \geq n$ 
if $\ell < n$, we may now put together \eqref{Upper bound Sigma1}, \eqref{eq:stalk} and \eqref{Equality Sigma3} to eventually obtain
\begin{align*}
\sum_{\nu = 1}^3 \Sigma^{(\nu)}(d,n;\ell) \leq~& M_1(d,n;\ell)^2
\left(1+ \boldsymbol{1}_{\ell=n}
+O\left(\left(\frac{3}{4}\right)^{\min\{\ell,\mu_\ell\}} +\frac{1}{\min\{d,n\}}
\right)\right).
\end{align*}
Recalling the equality \eqref{Equality M2} we see that this completes the proof. 
\end{proof}

\end{document}
